\documentclass[10pt]{article}

\usepackage[utf8]{inputenc}

\usepackage[T1]{fontenc}

\usepackage{amsmath}

\usepackage{amsfonts}

\usepackage{amssymb}

\usepackage[version=4]{mhchem}

\usepackage{stmaryrd}

\usepackage{bbold}

\usepackage{mathrsfs}



\begin{document}

Здесь $\sigma_{a}-$ Паули матрицы $(a=0,1,2,3), \sigma_{0}=I, \sigma_{ \pm}=$ $=(1 / 2)\left(\sigma_{1} \pm i \sigma_{2}\right), \bar{\psi}-$ функция, комплексно сопряженная функции $\psi, \varepsilon=1$ при $x>0, \varepsilon=i$ при $x<0$. Выполнение условия совместности для вспомогательной линейной задачи



\[

\frac{\partial U}{\partial t}-\frac{\partial V}{\partial x}+[U, V]=0

\]



эквивалентно выполнению уравнения (1). Запись уравнения (1) в виде (1') принято называть пре д с т а в ле н и е м нулевой кривизны.



Альтернативно метод обратной задачи рассеяния может быть сформулирован на основе Лакса представления.



Центральным объектом в методе обратной задачи рассеяния является матрица монодромии $T(\lambda)$. Для определения последней необходимо ввести матрицу перехода $T(x, y, \lambda)$, удовлетворяющую уравнению



\[

\frac{\partial}{\partial x} T(x, y, \lambda)=U(x, \lambda) T(x, y, \lambda)

\]



и условию



\[

\left.T(x, y, \lambda)\right|_{x=y}=I .

\]



Конкретное выражение матрицы монодромии через матрицу перехода зависит от вида граничных условий, накладываемых на функцию $\psi(x, t)$. Предположим, что решение Н.у.Ш. ищется в классе быстроубывающих функций с начальным условием $\psi(x) \in S\left(\mathbb{R}^{1}\right)\left[S\left(\mathbb{R}^{1}\right)-\right.$ пространство Шварца]. Тогда



\[

\begin{gathered}

T(\lambda)=\left\|\begin{array}{cc}

a(\lambda) & \varepsilon^{2} \bar{b}(\lambda) \\

b(\lambda) & \bar{a}(\lambda)

\end{array}\right\|=\lim e^{i \lambda x \sigma_{3} / 2} T(x, y, \lambda) e^{-i \lambda y \sigma_{3} / 2}, \\

x \rightarrow+\infty, \quad y \rightarrow-\infty .

\end{gathered}

\]



Замечательным свойством матрицы монодромии является особенно простая зависимость ее матричных элементов от времени:



\[

\begin{gathered}

a(\lambda, t)=a(\lambda, 0), \\

b(\lambda, t)=e^{-i \lambda^{2} t} b(\lambda, 0) .

\end{gathered}

\]



Функции $a(\lambda)$ и $b(\lambda)$ принято называть ко э фф и ц и е н там и перехода. В теории рассеяния величины $a^{-1}(\lambda)$ и $b(\lambda) / a(\lambda)$ играют роль коэффициентов прохождения и отражения. Решение $\psi(x, t)$ уравнения (1) однозначно восстанавливается по данным рассеяния и сводится к исследованию аналитич. свойств коэффициентов перехода. В частном случае безотражательного потенциала $[x<0, b(\lambda)=0]$ решение находится явно и называется $N$-с олитонны [где $N-$ число нулей коэффициентов $a(\lambda)$ ].



С помощью метода обратной задачи рассеяния также находится решение задачи Коши для граничных условий вида $\psi(x, t) \rightarrow \rho e^{i \varphi(t)}, x \rightarrow \pm \infty$ (условия конечной плотности). В этом случае обычно в правую часть уравнения (1) добавляют линейное по $\psi$ слагаемое $-2 x \rho^{2} \psi$ (соответственно в представлении нулевой кривизны матрица $V$ заменяется на $V=-i x \rho^{2} \sigma_{3}$ ).



В случае периодич. граничных условий $\psi(x+L, t)=$ $=\psi(x-L, t),-L \leqslant x \leqslant L$, решение Н.У.Ш. сводится к исследованию вспомогательной линейной задачи на римановой поверхности функции



\[

y^{2}=\prod_{n=1}^{\infty}\left(1-\frac{\lambda}{E_{n}}\right)

\]



Здесь $E_{n}-$ границы разрешенных и запрещенных зон в спектре оператора



\[

\mathscr{L}=i \sigma_{3} \frac{\partial}{\partial x}+i \sqrt{|x|}\left(\psi \sigma_{-}-\bar{\psi} \sigma_{+}\right) .

\]



В случае, когда число зон конечно, решение $\psi(x, t)$ уравнения (1) допускает явное выражение через $\theta$-функции Рима- на и называется конечнозон ным. При $L \rightarrow \infty$ конечнозонные решения Н.у. Ш. переходят в $N$-солитонные.



Н.у.Ш. можно рассматривать как гамильтоново уравнение



с гамильтонианом



\[

\frac{\partial \psi}{\partial t}=\{H, \psi\}

\]



\[

H=\int d x\left(\left|\frac{\partial \psi}{\partial x}\right|^{2}+x|\bar{\psi}|^{4}\right)

\]



и скобкой Пуассона $\{\psi(x), \bar{\psi}(y)\}=i \delta(x-y)$. Координатами в фазовом пространстве являются функции $\psi(x)$ и $\bar{\psi}(x)$ с определенными граничными условиями.



В рамках гамильтонова подхода к Н.у. Ш. широкое распространение получил метод $r$-матрицы, первоначально возникший в теории квантового метода обратной задачи. В основе данного метода лежит возможность представить скобки Пуассона матричных элементов матрицы $U(x, y)$ в виде:



\[

\begin{gathered}

\{U(x, \lambda) \otimes, U(y, \mu)\}= \\

=[r(\lambda-\mu), U(x, \lambda) \otimes I+I \otimes U(x, \mu)] \delta(x-y),

\end{gathered}

\]



где $r$-матрица равна



\[

r(\lambda-\mu)=-\frac{x}{2(\lambda-\mu)} \sum_{a=0}^{3} \sigma_{a} \otimes \sigma_{a}

\]



Можно показать, что такая запись скобок Пуассона заменяет представление нулевой кривизны.



Скобки Пуассона матричных элементов матрицы монодромии также записываются с помощью $r$-матрицы:



\[

\{T(\lambda) \otimes T(\mu)\}=[r(\lambda-\mu), T(\lambda) \otimes T(\mu)] .

\]



С точки зрения гамильтонова подхода переход к данным рассеяния является канонич. преобразованием к переменным действие-угол.



Гамильтонова модель Н. у. Ш. является вполне интегрируемой и обладает бесконечным набором интегралов движения $J_{n}$, производящей функцией для К-рых является след матрицы монодромии $\operatorname{tr} T(\lambda)$. Все интегралы движения записываются в виде локальных функционалов от $\psi$ и $\bar{\psi}$ и их производных, напр.:



\[

\begin{aligned}

& J_{1}=\int d x|\psi|^{2} \text { (оператор заряда), } \\

& J_{2}=\frac{i}{2} \int d x\left(\psi \frac{\partial \bar{\psi}}{\partial x}-\bar{\psi} \frac{\partial \psi}{\partial x}\right) \text { (оператор импульса), } \\

& J_{3}=H .

\end{aligned}

\]



Уравнения вида



\[

\frac{\partial \psi}{\partial t}=\left\{J_{n}, \psi\right\}

\]



принято называть в ы сш и м и Н.у. Ш.



В е к т о р о е Н.у. Ш. описывает движение заряженного скалярного поля $\psi_{a}(x, t), a=1, \ldots n$, с $n$ цветами:



\[

i \frac{\partial \psi}{\partial t}=-\frac{\partial^{2} \psi}{\partial x^{2}}+2 x \psi\left(\psi \psi^{*}\right)

\]



где под $\psi^{*}$ подразумевается вектор-столбец, эрмитово сопряженный к вектор-строке $\psi$. Векторное Н.у. Ш., как и скалярное, представимо в виде условия нулевой кривизны. Матрицы $U$ и $V$ в этом случае имеют размерность $(n+1) \times(n+1)$ и в блочной записи являются прямым обобщением матриц $U$ и $V$ скалярного уравнения. Гамильтониан модели и скобки Пуассона даются формулами



\[

\begin{aligned}

& H=\int d x\left[\frac{\partial \psi}{\partial x} \frac{\partial \psi^{*}}{\partial x}+x\left(\psi \psi^{*}\right)^{2}\right] \\

& \left\{\psi_{a}(x), \bar{\psi}_{b}(y)\right\}=i \delta_{a b} \delta(x-y)

\end{aligned}

\]



Иногда в литературе под термином «Н. у. Ш.» подразумевают систему уравнений



\[

\left.\begin{array}{l}

i \frac{\partial q}{\partial t}=-\frac{\partial^{2} q}{\partial x^{2}}+2 x q^{2} r \\

-i \frac{\partial r}{\partial t}=-\frac{\partial^{2} r}{\partial x^{2}}+2 x r^{2} q

\end{array}\right\}

\]





\end{document}